% !TEX root = ../jlt.tex
\section{The linter on the library}
\label{app:lint}

Table~\ref{tab:lint} counts what the linter of Section~\ref{sec:check} finds on the library behind the listings.
A check is a \emph{gate} on a file when the build fails unless the file has exactly the number of findings the repository states for it, most often none, and it is \emph{reported} on a file that is not yet clean: its findings are counted and printed, and the build does not fail.
Each check has a fixture that proves it fires, so a zero is not vacuous.
The library's module units~\cite[{\texttt{module.unit}}]{cpp23standard} are all in scope but the few interfaces that only re-export the others.
What each check matches, and how the tool runs, is the linter's reference, generated from its source (Appendix~\ref{app:derived}).

\makeatletter\let\lintinput\@@input\makeatother
\begin{table}[htbp]
\centering
\small
\begin{tabular}{@{}l c rrrr@{}}
\toprule
check & grade & gated & reported & with findings & findings \\
\midrule
\lintinput generated/lint-table.tex 
\bottomrule
\end{tabular}
\caption{The checks of \texttt{jlt-lint} on the library's module units in scope: the files on which each is a gate, an exact count, zero or not, that no change can move unseen; the files on which it is only reported; how many of all those have findings; and how many findings there are.
Each check is the file \texttt{tools/jlt-lint/rules/<check>.query}, a pattern over Clang's syntax tree, or a check of the linter that the file names.
Its grade says what its answer is worth, as a test of whether a definition shows, as written, the condition the check names: $\Pi$, rejecting, a finding refutes it and silence affirms nothing; $\Sigma$, affirming, an admission affirms it and its absence refutes nothing; $\Delta$, both; --- not a test.
For \texttt{admitted}, the last two columns count admissions.
The table is regenerated from the report of every build and compared with the one printed here.}
\label{tab:lint}
\end{table}

Three checks carry most of what the deny-rules find: \texttt{rule2-constrained-templates}, which counts templates with a parameter left unconstrained and is reported everywhere, \texttt{rule2-definition}, the definition check by names, and \texttt{rule1-local-state}.

\subsection{Admission by type}
\label{app:lint:admit}

The second stage of Section~\ref{sec:check} gives every type parameter of a constrained template an archetype, by reading the shapes that constrain it, with one rule for each kind of requirement a shape may have (Table~\ref{tab:archetype-rules}).
A shape with a requirement of any other kind leaves the templates written against it undecided, with that requirement named.
In the terms of Appendix~\ref{app:fjlt:typing}, a generic body reaches its shapes only as assumptions, a law by \textsc{E-Law-Assumed} and a required call by \textsc{T-Req-Assumed}; at the archetypes the same shapes are entailed instead, by \textsc{E-Shape}, \textsc{E-Law} and \textsc{E-Op}, from the entries the linter writes.
An archetype is the concrete counterpart of an assumed shape's record variable, the dictionary of Appendix~\ref{app:fjlt:translation}.
The standard's foundational concepts follow the Palo Alto report~\cite[\S 3.3]{stroustrup2012concept}: copying is the kernel's, equality is a call, and a default value is partially formed, to be written before it is read~\cite[\S 1.5]{stepanov2009elements}, the sugar for local state (Table~\ref{tab:derived}).
A parameter whose type is a constrained class template brings that class's constraints with it, since whoever calls the template already holds a value of it, as Rust's implied bounds propose~\cite{rust2017impliedbounds}.
A result type brings nothing, since \Cpp{} checks its constraints where the template is called, so a template whose result asks more than its constraints is refused: a caller who meets them could not call it.

\begin{table}[htbp]
\centering
\small
\begin{tabularx}{\linewidth}{@{}>{\raggedright\arraybackslash}p{0.34\linewidth}>{\raggedright\arraybackslash}p{0.3\linewidth}>{\raggedright\arraybackslash}X@{}}
\toprule
requirement & rules & archetype \\
\midrule
\multicolumn{3}{@{}l}{\emph{Kernel: the requirements $A$ of a shape (Figure~\ref{fig:fjlt-syntax})}} \\
$Q$: \texttt{copyable<T>}, \texttt{equality\_comparable<T>} & \textsc{E-Assume}; \texttt{==} by \textsc{T-Req-Assumed} / --- & \texttt{carrier}, \texttt{comparable\_carrier} \\
$Q$: \texttt{same\_as<T, U>} & solved / \textsc{E-Same} & unnamed \\
$e$: \texttt{unital<E>} & \textsc{E-Law-Assumed} / \textsc{E-Law} & unnamed \\
call: \texttt{E\{\}()} & \textsc{T-Req-Assumed} / \textsc{E-Op} & \texttt{nullary<C>} \\
call: \texttt{op\_t<E>\{\}(a, b)} & \textsc{T-Req-Assumed} / \textsc{E-Op} & \texttt{binary<C>} \\
call: \texttt{f(x)} & \textsc{T-Req-Assumed} / \textsc{E-Op} & \texttt{arrow<A, B>} \\
call: \texttt{f<T>(x)} & \textsc{T-Req-Assumed} / \textsc{E-Op}, \textsc{T-Call} & none \\
\midrule
\multicolumn{3}{@{}l}{\emph{Jlt: sugar (Table~\ref{tab:derived}) and the standard's concepts}} \\
\texttt{sized\_range<R>} & range; \textsc{E-Range} & \texttt{sized\_range\_of<C>} \\
\texttt{-> convertible\_to<T>} & \texttt{convertible\_to}; \textsc{T-Conv} & \texttt{converts\_to<T>} \\
\texttt{movable<T>}, \texttt{destructible<T>} & Rule~1: copyable data & \texttt{carrier} \\
\texttt{semiregular<T>}, \texttt{regular<T>} & local state; N3351 & \texttt{semiregular\_carrier}, \texttt{regular\_carrier} \\
\texttt{integral<N>}, \texttt{floating\_point<N>} & atoms $b$ & \texttt{every\_arithmetic} \\
\texttt{F::Domain}, \texttt{Presents<F::Codomain>} & member call; as \texttt{convertible\_to} & unnamed, its result \texttt{converts\_to} \\
a parameter \texttt{const Compose<F, G>\&} & RFC~2089 & unnamed \\
\bottomrule
\end{tabularx}
\caption{The rules by which the linter builds archetypes, one for each kind of requirement in a shape, in two groups: the kernel's requirements, and those of Jlt.
For a requirement of the kernel the middle column names the rule of Appendices~\ref{app:fjlt:typing} and~\ref{app:fjlt:translation} by which a generic body reaches it at its type variables, and the rule by which it is entailed at the archetypes; for one of Jlt it names the row of Table~\ref{tab:derived} or the source it follows.
An archetype is named as in the linter's header, \texttt{tools/jlt-lint/archetypes.hpp}, namespace \texttt{jlt::archetype}; ``unnamed'' marks a rule of the generator that builds no named type, and --- a requirement that Appendix~\ref{app:fjlt:typing} leaves to Jlt.}
\label{tab:archetype-rules}
\end{table}

\paragraph{The levels, at scale.}
\label{app:lint:summary}
Table~\ref{tab:summary} counts definitions, as Table~\ref{tab:levels} does for the listings, and by the same rule.
A definition is a declaration directly in the namespace judged, and the linter's list classes every kind of syntax it uses as kernel, sugar, scaffolding (modules, namespaces, assertions), refused, or still open.
It is refused when a deny-rule finds something in its lines, when it uses a kind of syntax the linter's list refuses, or when the type check does not admit it; undecided when it uses a kind still open or the type check leaves it undecided; in Jlt when every kind is kernel, scaffolding or sugar, and some kind is sugar; and in FJlt otherwise.
A kind is coarser than the grammar, so a level counts what a definition writes, not everything it means: the deny-rules and the type check supply the rest.
The table's last column counts definitions as well, each once, by the first outcome the type check gives within its lines.
On the library the type check refuses none, every definition it once refused having been fixed by stating the missing constraint, and what it leaves undecided is its generator's coverage, not findings: parameters of templates over values, requirements of forms it does not read, disjunctions and, by design, every body with an \cppinline{if constexpr}, of whose branches the archetypes check only one.
Every count is pinned per file and held by a ratchet, under which no body's count of definitions shown can fall without a reason written in the change.
The listings are the namespace of Listings~\ref{lst:motivation} to~\ref{lst:theorems}; the library is its module units in scope; Catch2 is its sources and its headers, each header judged as a file of its own, and the table reports the two apart.

\paragraph{The conjecture and its hypotheses.}
\label{app:lint:conjecture}
The conjecture of Section~\ref{sec:check} is that a definition admitted at its archetypes is well typed at every type that satisfies its shapes and keeps the rules itself, and resolves every choice its signature makes the same way.
Eight hypotheses exclude the counterexamples known to us.
\begin{enumerate}
\item A value is copyable data and a type's operations take their arguments by value or by constant reference, as Rule~1 asks, so that no argument the archetype accepts is refused by the type, and a shape that asks only \cppinline{movable} or \cppinline{destructible} is met by a copyable carrier.
\item A call by name inside a shape is qualified, so that no lookup by the argument's namespace, which \Cpp{20} extends to calls with template arguments, finds another function at the type.
\item The body does not depend on two carriers being different, or a carrier being different from an atom, which the archetypes make them.
\item A choice the signature leaves to the caller, as \cppinline{named_power} leaves the case of \cppinline{power}, may resolve differently.
\item The body observes no type except through a shape: a trait of the program read at a type variable outside a shape would be answered by the archetypes for themselves, and the standard's observations, \cppinline{sizeof}, \cppinline{alignof}, \cppinline{typeid}, its type traits and its concepts used as values, a deny-rule refuses.
\item An object is one value type: one presented by many types, as a family of expressions presents a set, need not convert to its codomain, which the archetypes do.
\item An entry stands above every use that reads it, which the linter checks only in part (Section~\ref{sec:limits}).
\item No limit of the implementation is reached.
\end{enumerate}
The linter excludes two more itself, a negation and a late conversion: at an archetype every law it does not declare is false, so a body whose callee asks \cppinline{!C<E>} is left undecided, and a result only convertible to its type is admitted only where it initialises a declared one.
In $F^G$, where a name is looked up where the template is defined, the corresponding statement is a theorem~\cite{siek2005essential}; Järvi et al.\ show what that binding costs the dispatch \Cpp{} programmers rely on~\cite{jarvi2006specialization}.
In \Cpp{} it is a conjecture, tested on the definitions of the listings and on the library.
